% ---
% id: reductions/wpfss/rem:tworoutes
% kind: rem
% title: "Soundness of the two routes and one route per artefact"
% summary: >-
%   Both routes are sound, and the encoded one, which needs no attestation, is
%   preferable where it is available. Pairing the licensed entry with the
%   encoded route's channel is unsound, so an artefact takes one route.
% prominence: standard
% status: accepted
% uses: [def:routes, reductions/cor:monotone, reductions/rem:channel-choice, prop:chain, def:targets, reductions/thm:transfer, reductions/rem:unsound]
% mentions: [reductions/thm:licence, reductions/def:link, reductions/prop:obstruction, req:one-route-per-artifact, reductions/req:constraint-and-licence-slots]
% assumes: []
% refs: []
% provenance: newly-derived
% lean:
% ---

Part~(iii) of Proposition~\ref{prop:chain} is the whole difference between the
two targets. Under $\Fwp$ the device fact is an operator inequality \emph{on
the target}, true of every state the adversary could possibly prepare, so it
needs no constraint, no licence and no attestation. Under $\Ffss$ the same
fact is true only on the image, so it must travel as a declared constraint
entry, and something must certify that the adversary cannot leave the image,
which is exactly Theorem~\ref{thm:licence}. Part~(ii) of the proposition says
which of the two hands the adversary more: by Corollary~\ref{cor:monotone} a
cleaner target never helps, and $\Ffss$ is cleaner than $\Fwp$, in general
strictly so: part~(iv) gives data on which the reverse relation fails.

\begin{remark}\label{rem:tworoutes}
Both routes of Definition~\ref{def:routes} are sound, and an artefact takes
exactly one of them (Requirement~\ref{req:one-route-per-artifact}). The encoded route is
preferable wherever it is available, because it requires no attestation a
program cannot check. It is \emph{not} always available:
Proposition~\ref{prop:obstruction} exhibits a device fact, a bound on a
\emph{union} of outcomes that must not be merged, that no target classical on
the whole of $\Qsys$ can carry and that $\Fwp$, one weight per outcome, does
not carry, and for which the licensed route is the only route; this is why
every artefact carries a constraint slot and a licence slot
(Requirement~\ref{req:constraint-and-licence-slots}).

\emph{Why one route per artefact.} The two routes run on different channels,
$\Lwp$ and $\Lfss$, and by Remark~\ref{rem:channel-choice}\,(c) a constraint
sound for one channel need not be sound for another. Here it is not: pairing
the licensed entry with the encoded route's channel gives, by
\eqref{eq:WPadjoint},
\begin{equation}\label{eq:crosspairing}
  \Lwp^{\dagger}\big(\Ffss_o-p_o\Pflag\big)
  =\blk{M}{\leqN}_o\ \dsum\
   \Big(\tfrac{1}{1-p}\big(\blk{M}{\gtN}_o-p_o\idon{}\big)-p_o\idon{}\Big),
\end{equation}
whose flag block is \emph{not} implied positive by \eqref{eq:certweights}: at
the tight end $\blk{M}{\gtN}_o=p_o\idon{}$ it equals $-p_o\idon{}\prec0$ for
$p_o>0$. So that pairing would impose a constraint some reachable state
violates, voiding Theorem~\ref{thm:transfer}, which is exactly the failure
Remark~\ref{rem:unsound} describes. Since an artefact's payload is
single-valued (Requirement~\ref{req:one-route-per-artifact}) and
\eqref{eq:link} binds each link's constraint list to \emph{that link's}
\tref{channel witness}, the remedy is structural: one route per artefact, and
two requests for a caller who wants both bundles, with no schema in which the
unsound pairing can be expressed.
\end{remark}
